\section{Честный статус и верификация}\label{sec:status}

Эта статья --- отражение построения на Lean~4, и её честность держится не на слове автора, а на
механическом аудите, который запускается при каждой сборке. Мы собираем здесь глобальный статус в
одном месте: что стоит машинно-доказанным при стандартных аксиомах, что условно на единственной
первопричине, что действительно открыто и как читатель может заново установить всё это из
исходного кода~\cite{lean4,mathlib,euclidspath-repo}.

\paragraph{Одна аксиома, один \code{sorry}.}
Всё построение опирается ровно на одну нестандартную аксиому --- \axiom{}, \emph{первопричину} из
\cref{sec:firstcause}, запертую в единственном карантинном модуле. Кроме неё используются лишь
собственные аксиомы Lean: \code{propext}, \code{Classical.choice} и \code{Quot.sound}. Точно так же
во всём корпусе ровно один \code{sorry}, и это сама цель --- \code{twin\_prime\_conjecture}. Ничего
иного не оставлено незавершённым; всякое иное утверждение либо \green{}~машинно-доказано при
стандартных аксиомах (включая зелёные условные теоремы, чьи гипотезы --- названные входы, а не
аксиома), либо \yellow{}~заражено аксиомой, то есть зависит от первопричины, либо \red{}~открыто
объявленный открытый вход. Зелёный модуль подделать нельзя, а \code{sorry} нельзя спрятать за
зелёной вывеской.

\paragraph{Перепись заражения.}
На момент письма верификатор сообщает \yellow{}~ровно $47$ заражённых аксиомой деклараций и
пересчитывает каждую при каждой сборке. Они распадаются на три группы. Сорок три живут в карантине:
аксиома \axiom{} и её теоремная свита в \srcpath{Engine/CausalClosureAxiom.lean} ($42$ декларации)
вместе с близнецовым следствием \code{higherEnergyIncompatibility\_twins}. Ещё две --- декларации
геометрии из \cref{sec:geometry}, наследующие заражение через близнецовую границу: встреча прямых,
которую нельзя узнать изнутри (\code{lines\_meet\_but\_unknowable\_from\_inside}), и близнецовые
двойники за всяким горизонтом (\code{twin\_vertices\_beyond\_every\_horizon}). Последние две ---
диадический слой начала, входящий через границу Навье--Стокса, а не через близнецовую: непричинимый
насос донной оболочки и его прочтение \code{dyadicBlowup\_is\_firstCauseManifestation}. Счётчик
заражения печатается, но не закрепляется; он честно растёт по мере добавления содержания, и вечными
должны оставаться лишь два инварианта ниже.

\paragraph{Верификатор и два его вечных инварианта.}
Скрипт \srcpath{scripts/VerifyAll.lean} перечисляет каждую невнутреннюю теорему и определение
\code{EuclidsPath} в собранном окружении, собирает через \code{collectAxioms} аксиомы, от которых
каждая транзитивно зависит, и отделяет три стандартные от всего прочего. Затем он
\emph{самоутверждает} два инварианта и валит запуск с ненулевым кодом выхода, если нарушен любой из
них:
\begin{equation}\label{eq:invariants}
\begin{gathered}
\text{(1) единственная нестандартная аксиома --- }\axiom{};\\
\text{(2) единственная заражённая \code{sorry} декларация --- }\code{twin\_prime\_conjecture}.
\end{gathered}
\end{equation}
Оба инварианта структурны, а не численны: они говорят, \emph{какая} аксиома и \emph{какой}
\code{sorry} могут существовать, но никогда --- сколько заражённых деклараций. Если бы в дерево
проскользнула посторонняя нестандартная аксиома или второе незавершённое доказательство,
верификатор отказался бы элаборировать. Именно это позволяет нам делать утверждения этой статьи, не
прося доверия.

\paragraph{Четыре состязательные вакуумности.}
Пустая теорема --- та, чьи гипотезы никогда не выполнимы, так что заключение ничего не говорит, ---
характерный провал формализованной витринной работы, и четырежды внутренний аудит поймал её раньше,
чем она попала в витрину. Мы их отмечаем, а не прячем. \emph{(i)} На близнецовой ветви вырожденный
пилинг $p = p\cdot 1$ сделал открытый узел опровержимым; заплата \code{properDiv} вместе с
\code{targetPos} вернула близнецовую гипотезу несущей, так что выполнимость на масштабе $A\ge 5$
действительно открыта. \emph{(ii)} На римановой ветви у маршрута со скачком ранга цель была не
заякорена: пакет \code{LiouvilleToTwinLocalization} населялся при нулевом входе
(\code{fullLocalization\_noInput}), а честной стеной оказалось ровно
$\neg\,\text{LiouvilleViolation}$ (\code{wall\_global}), силы гипотезы Римана. \emph{(iii)} На ветви
Мерсенна поздние пакеты \code{noEngine} были ненаселены, так что их заключения вакуумны и не могут
быть инстанцированы. \emph{(iv)} На ветви $\mathrm{P}$ против $\mathrm{NP}$ фронты извлечения,
запертые решателем, классически пусты (\code{PDecider} свободен), так что их выводы о разделении
вакуумны; классическое разделение \code{classicalSeparation\_of\_localIncompressible} остаётся
условным на запертом мосте и \emph{не} является доказательством $\mathrm{P}\ne\mathrm{NP}$. То, что
они пойманы машиной, а не рукой, --- сильнейший довод доверять частям, прошедшим витрину.

\paragraph{Ни одна задача тысячелетия не решена.}
Мы утверждаем это без оговорок. Ни одна из задач Клэя~\cite{millennium} здесь не решена, и о решении
ни одной не заявлено. \code{twin\_prime\_conjecture} остаётся \code{sorry}; гипотеза
Римана~\cite{clay-rh}, $\mathrm{P}$ против $\mathrm{NP}$~\cite{clay-pnp},
Янга--Миллса~\cite{clay-ym}, Навье--Стокса~\cite{clay-ns}, Ходжа~\cite{clay-hodge} и
Бёрча--Свиннертон-Дайера~\cite{clay-bsd} --- каждая прочитана через двигатель к \emph{структурной}
зелёной половине и \emph{аналитической} половине, которая либо принята единственной первопричиной
(\yellow{}), либо оставлена как названный открытый вход (\red{}). Редукции подлинны и машинно
проверены; редукция не есть доказательство.

\paragraph{Воспроизводимость.}
Всё вышесказанное можно вывести заново из открытого исходного кода. Из корня репозитория
\code{lake exe cache get} тянет кэш сборки \code{mathlib}, \code{lake build} проверяет всё
построение, а \srcpath{lake env lean scripts/VerifyAll.lean} заново устанавливает два инварианта
\cref{eq:invariants} и громко падает, если нарушен любой из них. Полный исходный код открыт по
адресу \url{https://github.com/elamaunt/Euclids-path}, и идентичный контент доступен как двуязычный
сайт по адресу \url{https://elamaunt.github.io/Euclids-path/}; оба несут один и тот же текст, на
английском и русском, вплоть до последней теоремы.
